\section{Related work}
\label{sec:related}

\subsection{Robustness of speech recognition under acoustic degradation}

The sensitivity of ASR to noise and channel distortions has been studied since
the earliest statistical recognisers~\cite{lippmann1997speech}. Surveys such
as that of Virtanen et al.~\cite{li2014overview} organise this
literature around feature enhancement, model adaptation, and multi-condition
training. Controlled evaluation infrastructure has been decisive for
progress. The Aurora framework~\cite{pearce2000aurora} standardised noise
types and signal-to-noise ratios for telephone-band recognition, and the CHiME
challenge series~\cite{barker2018fifth} extended such evaluation to
far-field, reverberant, and overlapping-speech conditions. Classical feature
transformations such as RASTA filtering~\cite{hermansky1994rasta} targeted the
same distortions from the signal-processing side, and synthetic room impulse
responses generated with the image method of Allen and
Berkley~\cite{allen1979image} and its modern
implementations~\cite{habets2006room} remain standard tools for controlled
reverberation studies. More recently, data augmentation strategies such as
SpecAugment~\cite{park2019specaugment} have embedded artificial degradation
into the training process itself, and weakly supervised training at scale, as
exemplified by Whisper~\cite{radford2023robust}, has produced recognisers
with markedly improved noise tolerance. The present work inherits the
methodological principle of these efforts, namely controlled and reproducible
degradation of the input signal, but applies it to black-box services rather
than to trainable models, which precludes training-based countermeasures and
makes the measurement itself the object of study.

\subsection{Cascaded and end-to-end speech translation}

The cascade versus end-to-end question has accompanied speech translation
since neural systems became practical. Early sequence-to-sequence systems
demonstrated that audio-to-text translation without an explicit recognition
stage is feasible~\cite{berard2016listen,weiss2017sequence}, and
Bentivogli et al.~\cite{bentivogli2021cascade} compared the two paradigms
systematically on IWSLT data, finding that end-to-end systems could match
cascades in translation quality while differing in error structure, with
cascades inheriting recognition errors and end-to-end systems producing more
fluent but occasionally freer output. The survey of Sperber and
Paulik~\cite{sperber2020speech} synthesises this literature, including work on
robustness of end-to-end systems to noisy input. Subsequent research
addressed noise robustness of end-to-end ST through self-training and data
selection~\cite{pino20_interspeech}. These studies evaluate models that the
researchers train themselves on public corpora. Robustness of the end-to-end component in
particular has been treated as a training-data problem, with noise-augmented
source audio, self-training on noisy transcripts, and domain adaptation as
the standard remedies. That framing presumes control over the model, which
is exactly what a service consumer lacks. When the translation component is
a fixed API, the robustness question shifts from how to train the model to
how to select among available models and architectures under measured
channel conditions, and that selection question is what the present
benchmark operationalises. The commercially relevant
question, namely how the two architectures behave when every stage is a
fixed, provider-operated service selected and combined by an integrator, has
received far less attention, and the present benchmark is designed
specifically for that setting.

\subsection{Evaluation of commercial speech services}

A third strand of related work evaluates speech services as they are
exposed to users. Benchmarks for large audio-language models, such as
AIR-Bench~\cite{yang2024airbench} for audio understanding and
VoiceBench~\cite{chen2024voicebench} for LLM-based voice assistants,
measure knowledge, comprehension, and instruction following, and they
establish the black-box, generative-evaluation methodology that the present
benchmark extends to the acoustic-robustness axis; neither, however,
characterises behaviour under controlled channel degradation, and both stop
at the response boundary rather than following errors into downstream
tasks. The propagation of recognition errors into translation has a long
history in the cascade literature, going back to the statistical framework
of Ruiz and Federico~\cite{ruiz2014assessing} for relating ASR error types
to downstream MT quality, but that line of work assumes access to the
models and corpora involved rather than opaque commercial endpoints.
Methodologically, black-box evaluations face constraints that mirror ours:
no access to model internals, no control
over model versions, and reliance on text normalisation to compare outputs
against references. They have, however, largely concentrated on clean or
moderately degraded audio, and on transcription accuracy alone. The broader
evaluation culture in speech and audio quality research, with its
calibrated degradation methodology and emphasis on reproducibility,
similarly targets the acoustic layer rather than the linguistic stages
downstream, and it is that tradition which the present service-level
benchmark follows, adding the service-level measurement that
intelligibility metrics cannot substitute for, as the dissociation
documented in Section~\ref{sec:results:snr} shows. The
contribution of the present study relative to this strand is twofold. First,
the evaluation is extended across the translation boundary, quantifying how
recognition differences propagate into translation quality for both cascade
and end-to-end architectures. Second, the acoustic degradations are applied
under full experimental control, with calibrated noise levels, a synthetic
room impulse response, codec and packet-loss emulation, and fixed random
seeds, which makes the entire measurement reproducible rather than dependent
on whatever happened to be in the test recordings.

\subsection{Positioning}

In summary, robustness benchmarks address open models but not commercial
services, service evaluations address commercial ASR but usually only in
clean conditions and only at the recognition boundary, and the
cascade-versus-end-to-end literature addresses research models rather than
provider-operated systems. CascadeBench combines the three perspectives. It
applies the controlled degradation methodology of the robustness literature
to the black-box commercial setting, extends the measurement through the
translation stage for both architectures, and adds a three-reference scoring
protocol, with one reference family-disjoint from every system under test,
that quantifies a bias source that any evaluation involving modern
generative reference translations must confront. A systematic search of the speech-translation and
audio-benchmarking literature, including the benchmarks cited above, found
no previous study that measures all of these dimensions simultaneously for
the current generation of speech APIs, and none that reports the linear
recognition-to-translation error pass-through regression that
Section~\ref{sec:results} documents; we state this as a finding of our
search rather than as a proof of absence.
